\chapter{Derived Categories.}
\VerdierPhysicalPageAnchor{0174}%
\VerdierSemanticAnchor{ch3d}
\begingroup
\small
\setlength{\parskip}{0pt}
\setlength{\emergencystretch}{1.5em}
\setlength{\abovedisplayskip}{4pt plus 2pt minus 2pt}%
\setlength{\belowdisplayskip}{4pt plus 2pt minus 2pt}%

\VerdierSemanticAnchor{ch3d-1}
\section{Definitions of Derived Categories.}

\noindent\textit{Throughout this section, unless otherwise stated,
$\mathcal A$ denotes an abelian category.}

\VerdierSemanticAnchor{ch3d-1-1}
\subsection{Notation and Terminology.}

\VerdierSemanticAnchor{ch3d-1-1-1}
\noindent\textbf{1.1.1.}\quad
Let $\mathcal A$ be an additive category (not necessarily abelian). We
denote by $\operatorname{comp}^{b}(\mathcal A)$ (respectively,
$\operatorname{comp}^{+}(\mathcal A)$, respectively,
$\operatorname{comp}^{-}(\mathcal A)$) the full subcategory of
$\operatorname{comp}(\mathcal A)$ defined by the complexes $X$ such that
$X^i=0$ for $|i|$ sufficiently large (respectively, $i$ sufficiently
small, respectively, $i$ sufficiently large). We denote by
$K^{b}(\mathcal A)$ (respectively, $K^{+}(\mathcal A)$, respectively,
$K^{-}(\mathcal A)$) the full subcategory of $K(\mathcal A)$ defined by the
objects of $\operatorname{comp}^{b}(\mathcal A)$ (respectively,
$\operatorname{comp}^{+}(\mathcal A)$, respectively,
$\operatorname{comp}^{-}(\mathcal A)$).

\VerdierSemanticAnchor{ch3d-1-1-2}
\noindent\textbf{1.1.2.}\quad
Now let $\mathcal A$ be an (abelian) category and let $*$ be one of the
typographical symbols $b$, $+$, $-$, or ``empty''. We denote by
$\operatorname{comp}^{*,b}(\mathcal A)$ (respectively,
$\operatorname{comp}^{*,+}(\mathcal A)$, respectively,
$\operatorname{comp}^{*,-}(\mathcal A)$) the full subcategory of
$\operatorname{comp}^{*}(\mathcal A)$ defined by the complexes $X$, objects
of $\operatorname{comp}^{*}(\mathcal A)$, such that $H^i(X)=0$ for $|i|$
large (respectively, for $i$ small, respectively, for $i$ large).
Similarly, we denote by $K^{*,b}(\mathcal A)$ (respectively,
$K^{*,+}(\mathcal A)$, respectively, $K^{*,-}(\mathcal A)$) the full
subcategory of $K^{*}(\mathcal A)$ defined by the complexes $X$, objects of
$K^{*}(\mathcal A)$, such that $H^i(X)=0$ for $|i|$ large (respectively,
$i$ small, respectively, $i$ large). Thus we obtain eight subcategories of
$K(\mathcal A)$:

\VerdierSemanticAnchor{ch3d-1-1-2-diagram}
\begin{center}
  \resizebox{0.88\textwidth}{!}{%
    $\begin{tikzcd}[ampersand replacement=\&,column sep=3.5em,row sep=1.7em]
      \& \& K^{+}(\mathcal A) \arrow[dr] \& \& \\
      \& K^{+,b}(\mathcal A) \arrow[ur] \arrow[dr]
        \& \& K^{,+}(\mathcal A) \arrow[dr] \& \\
      K^{b}(\mathcal A) \arrow[ur] \arrow[dr]
        \& \& K^{,b}(\mathcal A) \arrow[ur] \arrow[dr]
        \& \& K(\mathcal A) \\
      \& K^{-,b}(\mathcal A) \arrow[ur] \arrow[dr]
        \& \& K^{,-}(\mathcal A) \arrow[ur] \& \\
      \& \& K^{-}(\mathcal A) \arrow[ur] \& \&
    \end{tikzcd}$%
  }\hfill\textup{(1.1.2.1)}
\end{center}
\endgroup
